\section{Scope, Failure Modes, and Conclusion}
\label{sec:conclusion}

The method uses observational data through prediction, not through an assumption that the observational treatment contrast is causal. Trial randomization validates the pseudo-outcome for every fixed outcome regression. A common covariate marginal, or an exact density ratio under unequal marginals, centers the prediction-powered difference of covariate averages. These facts support separate \(\lambda\)- and \(\omega\)-channels for RCT-population CATE and ATE estimation.

The first boundary is covariate-marginal mismatch. Without weighting, if \(P_R(X)\neq P_O(X)\), the ATE correction has bias
\begin{align}
\omega
\left[
\mathbb E_O\{g(X)\}
-
\mathbb E_R\{g(X)\}
\right],
\end{align}
and the observational and trial terms in the unweighted CATE loss no longer cancel. The exact-ratio ATE theorem addresses the first problem when \(P_R^X\ll P_O^X\) and the true ratio \(r=dP_R^X/dP_O^X\) is available. With an estimated ratio \(\widehat r\), the exact conditional drift is \(\omega\mathbb E_O[(\widehat r-r)g]\). Evaluation-only variance estimation is valid at first order only in the stated negligible-drift regime; otherwise the ratio estimator's influence-function contribution must be included. The present CATE theorem remains a common-marginal result.

The second failure mode is weak trial overlap or an incorrect propensity. The conditional-mean identity in Proposition~\ref{prop:pseudo-validity} uses the known randomized-trial propensity and its strict positivity. Extreme inverse-propensity factors can also make the required moments unstable even when the identity remains algebraically valid.

The third failure mode is weak or misleading observational prediction. Under the common-marginal model, the unrestricted oracle \(\omega\)-gain is controlled by the squared correlation between \(g(X)\) and \(\tau(X)\), together with the finite \(n/N\) attenuation. A weak correlation gives little between-\(X\) gain. The \(\lambda\)-channel can likewise be unhelpful when \(Z_1-Z_0\) is noisy or nearly degenerate. Population oracle coefficients reveal these limits but do not ensure that their sample estimates are stable.

The fourth failure mode is adaptive reuse of evaluation data. Estimating nuisance functions, tuning \((\lambda,\omega)\), and evaluating the selected candidate on the same outcomes invalidates the conditioning used by the exact results. The explicit score-level counterexamples in Appendix~\ref{app:proofs} show that such reuse can create bias even when every fixed coefficient is unbiased. Three-way splitting prevents direct leakage. Nested outer cross-fitting preserves foldwise and aggregate unbiasedness under fold honesty, but overlapping training complements invalidate a naive sum of fold variances. A pooled interval needs an additional stability or influence-function argument.

The fifth failure mode is an overly restrictive or highly penalized CATE learner. Population target invariance does not imply that a finite-sample restricted projection is accurate. The square-completed representation is a convex mixture only for \(0\leq\omega<1\), and positive semidefiniteness is guaranteed only for \(0\leq\omega\leq1\) in function values or a linear sieve. Nonlinear parameterizations and coefficients outside this interval require direct analysis of the original quadratic.

The sixth failure mode is insufficient tail control. ATE variance identities require the relevant pseudo-outcomes and predictions to be square-integrable. Exact unbiasedness of the honest sample-variance estimator does not by itself imply confidence-interval coverage; the Wald theorem additionally uses a conditional Lindeberg condition, a sufficient \((2+\delta)\)-moment bound, and a nondegenerate variance limit. Loss-variance tuning requires \(U_{t,\lambda}\) and \(G_t\) to be square-integrable, which imposes fourth-moment-type conditions on residuals. Heavy tails can make the oracle formulas statistically unstable even when their denominators are nonzero.

Within these boundaries, the framework gives a model-agnostic organization of prediction-powered RCT--OBS fusion. The randomized trial fixes the causal target. The observational outcome model may reduce within-\(X\) pseudo-outcome noise through \(\lambda\), and the observational covariate sample may reduce predictable between-\(X\) variation through \(\omega\). For ATE estimation, the exact population oracle weakly improves the variance objective whenever the trial-only coefficient is feasible, but this is not a finite-sample no-harm theorem for estimated coefficients. For CATE learning, the same ingredients produce a target-preserving population risk and a precisely qualified empirical quadratic.

Empirical comparisons with the trial-only fallback should report selection variability, denominator instability, covariate-marginal mismatch, and performance under restricted learners. The exact results should be read as conditional identities, population oracle comparisons, and design guidance rather than as a universal no-harm guarantee.
